% ===== PRÉAMBULE CANONIQUE — à coller VERBATIM en tête de CHAQUE .tex =====
\documentclass[11pt,a4paper]{article}
\usepackage[utf8]{inputenc}\usepackage[T1]{fontenc}\usepackage{lmodern}
\usepackage[french]{babel}
\usepackage{amsmath,amssymb,mathtools}\usepackage{icomma}
\usepackage[version=4]{mhchem}          % \ce{...} : équations & formules
\usepackage{siunitx}\sisetup{locale=FR} % \qty{}{}, \si{}, \ang{}
\usepackage{chemfig}                    % \chemfig{...} : structures développées
\usepackage{xcolor}\usepackage{colortbl}
\usepackage{tikz}\usepackage{pgfplots}\pgfplotsset{compat=1.18}
\usepackage[most]{tcolorbox}\usepackage{enumitem}\usepackage{array,booktabs}
\usepackage[a4paper,margin=2.2cm]{geometry}
\usetikzlibrary{arrows.meta,calc,angles,quotes,positioning,patterns,backgrounds,shapes.geometric}
\definecolor{primary}{RGB}{222,49,99}\definecolor{primarydark}{RGB}{150,22,55}
\definecolor{defcol}{RGB}{0,114,178}\definecolor{methcol}{RGB}{52,92,148}
\definecolor{excol}{RGB}{0,135,90}\definecolor{attncol}{RGB}{200,40,55}
\tcbset{boxbase/.style={enhanced,breakable,boxrule=0.9pt,arc=2.5pt,
  left=3mm,right=3mm,top=2.3mm,bottom=2.3mm,fonttitle=\bfseries,coltitle=white}}
% --- boîtes TUTORIEL (noms EXACTS, ne pas renommer) ---
\newtcolorbox{cle}[1][]{boxbase,colback=defcol!6,colframe=defcol,
  title={\textsc{À retenir}\ifx\relax#1\relax\relax\else\textmd{ : \itshape #1}\fi}}
\newtcolorbox{methode}[1][]{boxbase,colback=methcol!7,colframe=methcol,sharp corners,
  title={\textsc{Méthode}\ifx\relax#1\relax\relax\else\textmd{ --- \itshape #1}\fi}}
\newtcolorbox{exemplebox}[1][]{boxbase,colback=excol!5,colframe=excol,
  title={\textsc{Exemple}\ifx\relax#1\relax\relax\else\textmd{ : \itshape #1}\fi}}
\newtcolorbox{piege}[1][]{boxbase,colback=attncol!6,colframe=attncol,
  title={\textsc{Piège}\ifx\relax#1\relax\relax\else\textmd{ : \itshape #1}\fi}}
\newtcolorbox{remarque}[1][]{boxbase,colback=black!4,colframe=black!45,
  title={\textsc{Remarque}\ifx\relax#1\relax\relax\else\textmd{ : \itshape #1}\fi}}
% --- boîtes EXERCICE / CORRIGÉ (noms EXACTS, auto-comptées) ---
\newtcolorbox[auto counter]{exo}[1][]{boxbase,colback=primary!4,colframe=primary,
  title={\textsc{Exercice}~\thetcbcounter\ifx\relax#1\relax\relax\else\textmd{ --- \itshape #1}\fi}}
\newtcolorbox[auto counter]{corrige}[1][]{boxbase,colback=excol!4,colframe=excol,
  title={\textsc{Corrigé}~\thetcbcounter\ifx\relax#1\relax\relax\else\textmd{ --- \itshape #1}\fi}}
\newcommand{\R}{\mathbb{R}}\newcommand{\N}{\mathbb{N}}\newcommand{\Z}{\mathbb{Z}}\newcommand{\ds}{\displaystyle}
% (un \usepackage{fancyhdr}+en-tête est possible ; il est ignoré côté web)
% =========================================================================
\begin{document}
\sisetup{per-mode=symbol}

\begin{exo}[déduire l'ordre des $\mathrm{p}K_a$ à partir de réactions]
On sait que les trois réactions suivantes sont \textbf{déplacées vers la droite} ($K_{\text{équi}} > 1$) :
\begin{enumerate}[label=\arabic*)]
  \item \ce{HNO2 + CH3COO- -> NO2- + CH3COOH}
  \item \ce{CH3COOH + HCO3- -> CH3COO- + H2CO3}
  \item \ce{H2CO3 + NH3 -> HCO3- + NH4+}
\end{enumerate}
Sans connaître les valeurs, déduis-en l'\textbf{ordre croissant} des $\mathrm{p}K_a$ des acides \ce{HNO2},
\ce{CH3COOH}, \ce{H2CO3} et \ce{NH4+}.
\end{exo}

\begin{exo}[quelles réactions sont déplacées vers la droite ?]
On donne les $\mathrm{p}K_a$ : \ce{HClO2} ($2{,}0$), \ce{H2CO3} ($6{,}4$), \ce{HCN} ($9{,}2$), et
\ce{CH3COOH} ($4{,}8$). Parmi les réactions suivantes, lesquelles sont déplacées vers la droite ?
\begin{enumerate}[label=\arabic*)]
  \item \ce{CH3COOH + ClO2- <=> CH3COO- + HClO2}
  \item \ce{HClO2 + CH3COO- <=> ClO2- + CH3COOH}
  \item \ce{HCN + HCO3- <=> CN- + H2CO3}
  \item \ce{H2CO3 + CN- <=> HCO3- + HCN}
\end{enumerate}
\end{exo}

\begin{exo}[retrouver la réaction d'une constante donnée]
On donne les $\mathrm{p}K_a$ : \ce{CH3COOH} ($4{,}8$), \ce{H2CO3} ($6{,}4$), \ce{NH4+} ($9{,}2$),
\ce{HCO3-} ($10{,}4$). Pour \textbf{laquelle} des réactions ci-dessous la constante d'équilibre vaut-elle
$K_{\text{équi}} = 10^{-2{,}8}$ ?
\begin{enumerate}[label=\alph*)]
  \item \ce{CH3COOH + HCO3- <=> CH3COO- + H2CO3}
  \item \ce{NH4+ + HCO3- <=> NH3 + H2CO3}
  \item \ce{H2CO3 + NH3 <=> HCO3- + NH4+}
  \item \ce{CH3COOH + NH3 <=> CH3COO- + NH4+}
\end{enumerate}
\end{exo}

\begin{exo}[quelles bases réagissent complètement avec \ce{HF} ?]
On veut savoir quelles bases réagissent \textbf{complètement} ($\Delta\mathrm{p}K_a > 2$) avec l'acide
fluorhydrique \ce{HF} ($\mathrm{p}K_a = 3{,}2$). On donne le $\mathrm{p}K_a$ de l'acide conjugué de chaque base :
\ce{NO2-} (\ce{HNO2}, $3{,}3$), \ce{CH3COO-} (\ce{CH3COOH}, $4{,}8$), \ce{CN-} (\ce{HCN}, $9{,}2$),
\ce{CO3^{2-}} (\ce{HCO3-}, $10{,}4$), \ce{ClO2-} (\ce{HClO2}, $2{,}0$).
\begin{enumerate}[label=\alph*)]
  \item Pour chaque base, calcule $\Delta\mathrm{p}K_a$ de la réaction avec \ce{HF}.
  \item Quelles bases réagissent \textbf{complètement} avec \ce{HF} ?
\end{enumerate}
\end{exo}

\begin{exo}[capstone : prédominance d'un diacide et piège du pH]
L'acide oxalique \ce{H2C2O4} est un diacide : \ce{H2C2O4/HC2O4-} ($\mathrm{p}K_{a1} = 1{,}2$) et
\ce{HC2O4-/C2O4^{2-}} ($\mathrm{p}K_{a2} = 4{,}3$).
\begin{enumerate}[label=\alph*)]
  \item Quelle espèce prédomine à $\mathrm{pH} = 3{,}0$ ?
  \item Dans l'intestin, à $\mathrm{pH} = 6{,}4$, quelle espèce prédomine ? (Attention au piège : ce n'est pas
        \ce{HC2O4-}.)
  \item À $\mathrm{pH} = 4{,}3$, quel est le rapport $\dfrac{[\ce{C2O4^{2-}}]}{[\ce{HC2O4-}]}$ ?
\end{enumerate}
\end{exo}

\end{document}
